% GENERATED FILE. Edit canonical lessons, then run npm run generate:books.
% canonical-source-hash: fnv1a64:75d1bb9f5d15a890
% canonical-lessons: UR-C467-talu, UR-C467-masurha, UR-C467-chatna, UR-C467-chusna, UR-C467-thukna

\chapter{Taste and the Mouth}
\label{ch:ur-taste-and-the-mouth}
\hlchaptermodality{467}

\begin{chapteropening}
\textbf{By the end of this chapter:} \emph{I can talk about the mouth and tasting.}

Complete the last lesson of chapter 467: I can talk about the mouth and tasting.
\end{chapteropening}

\section[\texorpdfstring{\ur{تالو} (tālū)}{تالو (tālū)}]{\ur{تالو} (tālū) — the palate}

\label{lesson:UR-C467-talu}



\begin{quote}
Before the new one: say the Urdu for the sole of the foot, then the Urdu for a paw.
\end{quote}

\subsection*{You'll want to know: \ur{تالو}}
\textbf{\ur{تالو}} — \emph{tālū} — \textquotedblleft{}the palate\textquotedblright{}.

\begin{quote}
\ur{تالو} \ur{منہ} \ur{کی} \ur{چھت} \ur{ہے۔}
\end{quote}

\begin{quote}
\emph{tālū muṅh kī chhat hai.}
\end{quote}

\begin{quote}
The palate is the roof of the mouth.
\end{quote}

Masculine.

\begin{cousinweb}
\textbf{\ur{تالو}} (\emph{tālū}) is inherited from Sanskrit \emph{tālu}, \textquotedblleft{}palate\textquotedblright{} (CDIAL 5803).
\end{cousinweb}

\begin{grammarlens}[title={What you've built}]
One more word for describing what you lived and felt.
\end{grammarlens}

\subsection*{Your turn}
\begin{itemize}
\raggedright
  \item \emph{Say it:} \emph{tālū}
  \item \emph{Say it:} \emph{tālū}, once more
  \item \emph{Say it:} \emph{tālū muṅh kī chhat hai.} — The palate is the roof of the mouth.
  \item \emph{Recall:} say the Urdu for the sole of the foot, then the Urdu for a paw, then say \emph{tālū} again
\end{itemize}

\begin{tcolorbox}[breakable,colback=teal!4,colframe=teal!35!black,title={Before you move on}]
What does \ur{تالو} mean? (\textquotedblleft{}The palate\textquotedblright{}.) Say it once more.
\end{tcolorbox}
\section[\texorpdfstring{\ur{مسوڑھا} (masūṛhā)}{مسوڑھا (masūṛhā)}]{\ur{مسوڑھا} (masūṛhā) — a gum (of the teeth)}

\label{lesson:UR-C467-masurha}



\begin{quote}
Before the new one: say the Urdu for a paw, then the Urdu for the palate.
\end{quote}

\subsection*{You'll want to know: \ur{مسوڑھا}}
\textbf{\ur{مسوڑھا}} — \emph{masūṛhā} — \textquotedblleft{}a gum (of the teeth)\textquotedblright{}.

\begin{quote}
\ur{دانت} \ur{کے} \ur{درد} \ur{سے} \ur{مسوڑھا} \ur{سوج} \ur{گیا۔}
\end{quote}

\begin{quote}
\emph{dānt ke dard se masūṛhā sūj gayā.}
\end{quote}

\begin{quote}
The gum swelled up with the toothache.
\end{quote}

Masculine.

\begin{cousinweb}
\textbf{\ur{مسوڑھا}} (\emph{masūṛhā}) is inherited: CDIAL lists the Hindi cousin \emph{masūṛā} under \emph{māṃsapuṭaka}, \textquotedblleft{}the gums\textquotedblright{}, a form it reconstructs from the words for \textquotedblleft{}flesh\textquotedblright{} and \textquotedblleft{}fold\textquotedblright{} (CDIAL 9984).
\end{cousinweb}

\begin{grammarlens}[title={What you've built}]
One more word for describing what you lived and felt.
\end{grammarlens}

\subsection*{Your turn}
\begin{itemize}
\raggedright
  \item \emph{Say it:} \emph{masūṛhā}
  \item \emph{Say it:} \emph{masūṛhā}, once more
  \item \emph{Say it:} \emph{dānt ke dard se masūṛhā sūj gayā.} — The gum swelled up with the toothache.
  \item \emph{Recall:} say the Urdu for a paw, then the Urdu for the palate, then say \emph{masūṛhā} again
\end{itemize}

\begin{tcolorbox}[breakable,colback=teal!4,colframe=teal!35!black,title={Before you move on}]
What does \ur{مسوڑھا} mean? (\textquotedblleft{}A gum (of the teeth)\textquotedblright{}.) Say it once more.
\end{tcolorbox}
\section[\texorpdfstring{\ur{چاٹنا} (chāṭnā)}{چاٹنا (chāṭnā)}]{\ur{چاٹنا} (chāṭnā) — to lick}

\label{lesson:UR-C467-chatna}



\begin{quote}
Before the new one: say the Urdu for the palate, then the Urdu for a gum.
\end{quote}

\subsection*{You'll want to know: \ur{چاٹنا}}
\textbf{\ur{چاٹنا}} — \emph{chāṭnā} — \textquotedblleft{}to lick\textquotedblright{}.

\begin{quote}
\ur{بلی} \ur{نے} \ur{دودھ} \ur{چاٹا۔}
\end{quote}

\begin{quote}
\emph{billī ne dūdh chāṭā.}
\end{quote}

\begin{quote}
The cat licked up the milk.
\end{quote}

Past, with \ur{نے} (he or she …): \ur{اس} \ur{نے} \ur{چمچ} \ur{چاٹا}, \emph{us ne chammach chāṭā}.

\begin{cousinweb}
\textbf{\ur{چاٹنا}} (\emph{chāṭnā}) is inherited from \emph{caṭṭ}, \textquotedblleft{}lick, taste\textquotedblright{}, a form CDIAL reconstructs (CDIAL 4573).
\end{cousinweb}

\begin{grammarlens}[title={What you've built}]
One more word for describing what you lived and felt.
\end{grammarlens}

\subsection*{Your turn}
\begin{itemize}
\raggedright
  \item \emph{Say it:} \emph{chāṭnā}
  \item \emph{Say it:} \emph{chāṭnā}, once more
  \item \emph{Say it:} \emph{billī ne dūdh chāṭā.} — The cat licked up the milk.
  \item \emph{Recall:} say the Urdu for the palate, then the Urdu for a gum, then say \emph{chāṭnā} again
\end{itemize}

\begin{tcolorbox}[breakable,colback=teal!4,colframe=teal!35!black,title={Before you move on}]
What does \ur{چاٹنا} mean? (\textquotedblleft{}To lick\textquotedblright{}.) Say it once more.
\end{tcolorbox}
\section[\texorpdfstring{\ur{چوسنا} (chūsnā)}{چوسنا (chūsnā)}]{\ur{چوسنا} (chūsnā) — to suck}

\label{lesson:UR-C467-chusna}



\begin{quote}
Before the new one: say the Urdu for a gum, then the Urdu for to lick.
\end{quote}

\subsection*{You'll want to know: \ur{چوسنا}}
\textbf{\ur{چوسنا}} — \emph{chūsnā} — \textquotedblleft{}to suck\textquotedblright{}.

\begin{quote}
\ur{بچہ} \ur{اپنا} \ur{انگوٹھا} \ur{چوستا} \ur{ہے۔}
\end{quote}

\begin{quote}
\emph{bachchā apnā angūṭhā chūstā hai.}
\end{quote}

\begin{quote}
The baby sucks his thumb.
\end{quote}

Past, with \ur{نے} (he or she …): \ur{اس} \ur{نے} \ur{آم} \ur{چوسا}, \emph{us ne ām chūsā}.

\begin{cousinweb}
\textbf{\ur{چوسنا}} (\emph{chūsnā}) is inherited from Sanskrit \emph{cūṣati}, \textquotedblleft{}sucks\textquotedblright{} (CDIAL 4898).
\end{cousinweb}

\begin{grammarlens}[title={What you've built}]
One more word for describing what you lived and felt.
\end{grammarlens}

\subsection*{Your turn}
\begin{itemize}
\raggedright
  \item \emph{Say it:} \emph{chūsnā}
  \item \emph{Say it:} \emph{chūsnā}, once more
  \item \emph{Say it:} \emph{bachchā apnā angūṭhā chūstā hai.} — The baby sucks his thumb.
  \item \emph{Recall:} say the Urdu for a gum, then the Urdu for to lick, then say \emph{chūsnā} again
\end{itemize}

\begin{tcolorbox}[breakable,colback=teal!4,colframe=teal!35!black,title={Before you move on}]
What does \ur{چوسنا} mean? (\textquotedblleft{}To suck\textquotedblright{}.) Say it once more.
\end{tcolorbox}
\section[\texorpdfstring{\ur{تھوکنا} (thūknā)}{تھوکنا (thūknā)}]{\ur{تھوکنا} (thūknā) — to spit}

\label{lesson:UR-C467-thukna}



\begin{quote}
Before the new one: say the Urdu for to lick, then the Urdu for to suck.
\end{quote}

\subsection*{You'll want to know: \ur{تھوکنا}}
\textbf{\ur{تھوکنا}} — \emph{thūknā} — \textquotedblleft{}to spit\textquotedblright{}.

\begin{quote}
\ur{سڑک} \ur{پر} \ur{نہیں} \ur{تھوکنا} \ur{چاہیے۔}
\end{quote}

\begin{quote}
\emph{saṛak par nahīṅ thūknā chāhiye.}
\end{quote}

\begin{quote}
You should not spit in the street.
\end{quote}

Past, with \ur{نے} (he or she …): \ur{اس} \ur{نے} \ur{پانی} \ur{تھوکا}, \emph{us ne pānī thūkā}.

\begin{cousinweb}
\textbf{\ur{تھوکنا}} (\emph{thūknā}) is inherited from \emph{thukk}, a form CDIAL reconstructs and treats as an imitation of the sound (CDIAL 6097).
\end{cousinweb}

\begin{grammarlens}[title={What you've built}]
One more word for describing what you lived and felt.
\end{grammarlens}

\subsection*{Your turn}
\begin{itemize}
\raggedright
  \item \emph{Say it:} \emph{thūknā}
  \item \emph{Say it:} \emph{thūknā}, once more
  \item \emph{Say it:} \emph{saṛak par nahīṅ thūknā chāhiye.} — You should not spit in the street.
  \item \emph{Recall:} say the Urdu for to lick, then the Urdu for to suck, then say \emph{thūknā} again
\end{itemize}

\begin{tcolorbox}[breakable,colback=teal!4,colframe=teal!35!black,title={Before you move on}]
What does \ur{تھوکنا} mean? (\textquotedblleft{}To spit\textquotedblright{}.) Say it once more.
\end{tcolorbox}
